\chapter{Conclusion}
\label{conclusion}

Probabilistic models for football outcomes have been studied extensively, specifically for European competitions.
However, the validation via Simulation-Based Calibration and the extension to the player level remain uncommon.
This work addresses this gap by making three main contributions.

The first contribution is a dataset spanning from 2013 to 2025, covering five competitions (\textit{Série A} through \textit{D} and \textit{Copa do Brasil}).
It maps over 16,800 unique players across 368 clubs and is available at \url{https://github.com/BrazilianFootball/Data}, with an archived version on Zenodo \cite{michels_2026_dataset}.
The extraction pipeline converts PDF match reports into structured JSON files with match-level summary and temporal segmentation based on substitutions and red cards.
This information enables the analysis at both club and player levels.

The second contribution is the application of Simulation-Based Calibration to validate the inference procedure for all 22 proposed models.
The validation occurs for 14 club-level (4 Bradley-Terry and 10 Poisson variants) and 8 player-level (all Poisson variants) specifications.
All of them passed validation, with flagging rates consistent with chance under 95\% confidence bands.
This procedure ensures that posterior distributions are correctly calibrated before application to real data.

The third contribution is the comparative evaluation of models on \textit{Brasileirão} data from 2020--2025 and their extension to the player level.
Model P7 (single skill with home effect and common effect) achieved the best predictive performance according to Brier Score, Log Score, and Ranked Probability Score.
Analysing year by year, the home effect remained robust over time, with estimates for the multiplicative effect on the home team's expected goals ranging from 1.31 to 1.53.
The common effect parameter was statistically indistinguishable from zero, indicating that team strengths alone adequately capture the goal-scoring structure in \textit{Brasileirão}.

At the player level, log-skill parameters were estimated for 730 unique players in the single-season analysis and 1,975 in the multi-season analysis.
The synthetic player parameter quantifying red card effects showed strong negative impact on the team's expected goals when aggregating players to the team level.
For regular players, all estimated credible intervals include zero, indicating that the model fails to identify any significant difference between players.
Furthermore, individual skill estimates correlate near-perfectly with weighted goal difference.
This is a strong indication that the player-level model captures the team context in matches played by the player rather than isolated individual ability.

This confounding between player skill and team performance reflects a fundamental limitation of models relying solely on match outcomes.
Without individual action data, the contribution of each player cannot be separated from the collective team effect.
More broadly, the models assume that team strengths are constant within each season and that matches are exchangeable.
This assumption ignores several factors that may influence outcomes, such as fatigue from fixture congestion and the differential pressure of decisive versus inconsequential matches.
Under exchangeability, a win in the first round carries the same inferential weight as a win in the final round, precluding the detection of form fluctuations or momentum effects.

Several directions could address these limitations.
Incorporating player-specific event data would enable finer performance attribution and reduce the confounding between individual skill and team context.
Relaxing the exchangeability assumption would require dynamic specifications where team and player abilities evolve over time.
State-space formulations, in which parameters follow a random walk or autoregressive process, would allow the model to capture form fluctuations.
Such extensions would also enable predictions that account for recent performance rather than treating all historical matches equally.
Another possibility is adding contextual covariates, such as days since last match, travel distance, or an indicator for decisive matches, to the model.
These covariates would allow the model to distinguish between high- and low-pressure moments without abandoning the exchangeability framework entirely.

To support reproducibility, all code used in this work is available at \url{https://github.com/BrazilianFootball/Modeling}.
